\subsection{A verified finite input route and an indexed dual reduction}
\label{ssec:certified-input-route}

The next computation strengthens the insufficient-mode example to the
ground state of the finite analytic matrix. At $C=250,N=500$, the matrix
is assembled from digamma and trigamma formulas with finite endpoints and prime
powers, with Arb rounding and explicit geometric tail bounds. No retained
eigenstate or zeta ordinate enters this matrix/eigenpair construction.
The interval implementation encloses 1830 values from an independent
high-precision assembly control. These checks validate the implementation;
the analytic identities and remainder bounds remain its mathematical basis.

\paragraph{Enclosed scalar branch.}
At a real shift $z$, positive pivots of $B-zI$ certify its positivity.
The central Schur jet yields
\[
 \sigma(z)=a-z-c^T(B-zI)^{-1}c,\quad
 -\sigma'(z)=1+\eta(z),\quad
 b(z)=\frac1{\operatorname{tr}(B-zI)^{-1}}
       \le\lambda_{\min}(B-zI).
\]
Here $\eta(z)=\|(B-zI)^{-1}c\|^2$. If $0<a_*<b(z)$, then throughout
$|t|\le a_*$ the magnitude of the Schur derivative lies between
\[
 M_-=1+\frac{\eta(z)}{(1+a_*/b(z))^2},\qquad
 M_+=1+\frac{\eta(z)}{(1-a_*/b(z))^2}.
\]
Verifying $a_*M_->|\sigma(z)|$ gives opposite endpoint signs and the
enclosure $z+\sigma(z)/[M_-,M_+]$ for its unique root. The positive
noncentral block identifies this root as the lowest finite eigenvalue.

The direct interval elimination fails to prove positive pivots at 1000
and 3000 working digits, and succeeds at 5000. Seven scalar refinements
give relative eigenvalue radius below $7.47\cdot10^{-1415}$. Direct
interval vector back-substitution remains too wide. Instead, a point
central-normalized vector $v=(1,v_{\rm tail})$ is checked by recomputing
\[
 r_v=(B-\varepsilon I)v_{\rm tail}+c
\]
in interval arithmetic. A verified $g\le\lambda_{\min}(B-\varepsilon I)$
gives $\|v_*-v\|\le\|r_v\|/g=:E$. If $E<\|v\|$, the unit-vector
error is at most $2E/(\|v\|-E)$. The resulting bound is
$2.271\cdot10^{-1424}$. The large guard precision is a feature of this
certification algorithm, not a necessary physical mode or precision budget.

\paragraph{A finite-operator mode breaker.}
Five hundred disjoint positive interval-Newton root enclosures, together
with a nonzero leading coefficient of the degree-500 numerator, exhaust
its movable roots for the enclosed ground vector. Previous roots are warm
seeds only. The first 362 movable roots match their corresponding zeta
midpoint windows, and the $k=363$ window contains no movable root. Every
one of the 362 earlier displacement intervals contains the new bound.
Thus this breaker now concerns the independently enclosed finite analytic
operator. Carrier roots are still counted separately. The later regularization in
Section~\ref{ssec:ground-regularization} supplies an analytic certificate
for the $C=250,N=1500$ thousand-root ground source.

The historical 5047-bit unit vector differs from the new finite ground
state by approximately $1.36603\cdot10^{-858}$; its eigenvalue has relative
difference $6.51756\cdot10^{-849}$. These are finite source errors consistent
with the retained root measurements, not a storage-corruption finding.

\paragraph{A failed larger-cutoff extrapolation.}
The fixed global-germ rule from Section~\ref{sec:transfer-germ-bounds}
does not give uniform higher-cutoff edge accuracy. The following counts
use a one-percent \emph{inverse-squared-slope} gate, rather than the earlier
combined-predictor gate:
\[
\begin{array}{c|rr}
(C,N)&64\text{ coefficients, }[31/32]&128\text{ coefficients, }[63/64]\\
\hline
(1000,890)&33/58&34/58\\
(1200,970)&4/10&3/10
\end{array}
\]
These are all eligible previously acquired points, not complete certified
prefixes or blind data. Mismatched and unsuccessful acquisitions remain
recorded. The higher-order system is first unresolved at 1100 working
digits; all targets are scored at 1700 digits and the model failures
remain. More coefficients do not repair it uniformly, and some misses
occur at roots that still have several accurate digits.

\paragraph{Two global vectors determine the indexed dual.}
Write $B_\varepsilon=B-\varepsilon I$, $D=\operatorname{diag}(1^2,\ldots,N^2)$,
and the inherited displacement identity as
\[
 DB_\varepsilon-B_\varepsilon D=rw^T-wr^T,\qquad
 w_j=\sqrt2,\quad r_j=\sqrt2\,j b_j.
\]
Set $X=B_\varepsilon^{-1}$, $p=Xr$, $q=Xw$, and
$R_t=(D-tI)^{-1}$ away from the finite frequency poles. Then
\begin{equation}\label{eq:input-cauchy-dual}
 XR_tw=R_t\{q(1+p^TR_tw)-p(q^TR_tw)\}.
\end{equation}
Indeed, $DX-XD=qp^T-pq^T$ and
$XR_t-R_tX=R_t(DX-XD)R_t$. Applying this identity to $w$ proves the
formula. Once $p,q$ are known, each indexed dual needs two Cauchy sums and
linear work in $N$, rather than another dense inverse solve. Numerical
cancellation in this identity still requires control.

For $t=(z/\beta)^2$, put
$K(z)=2z\sin(zL/2)/(\sqrt L\,\beta^2)$. The tail transform functional
is $-K(z)R_tw$. Thus the exact central-normalized finite transform obeys
\begin{equation}\label{eq:input-residual-transform}
 F_*(z)-F_v(z)=K(z)(XR_tw)^Tr_v.
\end{equation}
This follows directly from
$(v_*-v)_{\rm tail}=-Xr_v$. Both the inverse-displacement identity and
this sign convention pass 156 exact rational component checks.
A derivative lower bound and a valid root neighborhood turn
\eqref{eq:input-residual-transform} into an indexed displacement bound.
The specific norm $\|XR_tw\|$ can be much smaller than the generic inverse
norm. Its uniform dependence on $C,N,k$, together with the structured
residual's dependence on $P$, is a concrete remaining target for the
convergence law. The finite computation and exact reduction do not yet
establish that rate.
